\scanpage{9}
\mspage{6}

J'indique tout de suite l'axiome le plus important sur un dérivateur (au delà de la donnée d'un 2-\unsure{contrafoncteur|contr. foncteur} strict) :

\begin{tabular}{||p{.88\linewidth}}
\uline{Der}$_{1}$ : Si $f:I\longrightarrow J$ dans $\mathrm{Dia}$, alors $\deleted{\unsure{\mathbb D f}}\inserted{f^{*}_{\unsure{\mathbb D|D}}=}f^{*}:\mathbb D(J)\longrightarrow\mathbb D(I)$ a un adjoint à droite et un adjoint à gauche.
\end{tabular}

Ces foncteurs seront notés (du moins provisoirement) $f_{*}^{\mathbb D}$ ou $f_{*}$, et $f_{!}^{\mathbb D}$ ou $f_{!}$,

\[
f_{!},f_{*}:\mathbb D(I)\rightrightarrows\mathbb D(J)\ ,
\]

de sorte qu'on trouve une suite de \uline{trois} foncteurs adjoints entre $\mathbb D(I)$ et $\mathbb D(J)$

\[
f_{!}, f^{*}, f_{*}
\]

\uline{N.B.} Pour \deleted{\unsure{pour}} la plupart des développements on n'aura pas besoin \inserted{de l'existence} à la fois des $f_{!}$ et des $f_{*}$. Un \uline{dérivateur à gauche} ne suppose que \inserted{les} $f_{*}$ (adjoints à droite \inserted{de $f^{*}$}), et l'existence de $f_{!}$ \unsure{seulemt|seulement} quand $f$ est une \uline{immersion ouverte}. Dualement, un \uline{dérivateur à droite} ne suppose que les $f_{!}$ (adjoints à gauche des $f^{*}$), et $f_{*}$ seulemt si $f$ est une \uline{imm. fermée}. Ce sont sans doute les dér. à g. qui sont les plus importants.

% archive sheet number: 8
